% Ad Atticum 4.10 --- headnote
% ad-atticum-04-10, 23 April 55 BC, written at Cumae

\textit{Cicero to Atticus, written at Cumae on the ninth day
before the Kalends of May, 55 BC --- 22 or 23 April. The
opening news comes from down the road at Puteoli: that
Ptolemy XII Auletes, expelled from Egypt in 58 BC and the
subject of all the previous year's senatorial fight in the
Lentulus Spinther correspondence, has at last been restored
--- as in fact he was, by Gabinius's army from Syria in
spring 55. The political concerns of Cicero's hosts here
in the bay are now those of Faustus Cornelius Sulla, in
whose library Cicero is reading.}

\textit{The famous central piece is the comparison: ``I
would rather sit in that little chair of yours that you have
under the bust of Aristotle than in the curule chair of
those gentlemen, and walk with you at your house than with
the man with whom I see I must be walking.'' The walking
companion is Pompey, who is on his way over from his own
Cumae estate after arriving on the Parilia (21 April). The
contrast is the durable Ciceronian formula --- the friend's
study against the second consul's curule throne, the
philosopher's company against the politician's --- and is
made the more pointed by the closing report that Cicero is
in fact going to walk with the man the next morning. ``About
that walking, fortune will see, or some god if any god takes
thought.''}

\textit{The closing paragraph drops back to the building works
on the Palatine: Atticus is to chase Philotimus on the
walkway, the Laconian sweat-room, and the works in the manner
of Cyrus the architect. The reckoning of household business
sits beside the philosophical preference for the friend's
study; the letter is short, and in the bay April light, it
is the very thing.}
